\documentclass[11pt,a4paper]{article}
\usepackage[english]{babel}
\usepackage[utf8]{inputenc}
\usepackage[T1]{fontenc}
\usepackage{geometry}
\geometry{margin=2.5cm}
\usepackage{enumitem}
\usepackage{amsmath,amssymb}
\usepackage{hyperref}
\usepackage{xcolor}

\newcommand{\major}[1]{\textbf{\color{red!70!black}M#1.}}
\newcommand{\minor}[1]{\textbf{\color{orange!80!black}m#1.}}

\title{Mock Referee Report --- JASA\\[6pt]
\large\textit{Permutations accelerate Approximate Bayesian Computation}}
\author{Anonymous Reviewer}
\date{April 2026}

\begin{document}
\maketitle

\noindent\textbf{Recommendation:} Major revision.

\section*{Summary}

The paper introduces \textsc{permABC}, an ABC framework for hierarchical models with exchangeable compartments. The core idea---solving a linear sum assignment problem to optimally match simulated compartments to observed ones before checking the ABC criterion---is elegant and practically relevant. The authors develop sequential extensions (\textsc{permABC-SMC}), as well as two novel bridging strategies (Over-Sampling and Under-Matching), and demonstrate the approach on a COVID-19 epidemic model with 94 French departments ($d=283$).

The contribution is original and the practical motivation is strong. However, the paper suffers from a gap between its ambitions and what is actually demonstrated. The theoretical results are limited in scope, the numerical evaluation relies almost exclusively on a single proxy metric (tolerance~$\varepsilon$), and several claims lack adequate empirical or theoretical support.

\section*{Major concerns}

\major{1} \textbf{The sole evaluation metric is $\varepsilon$---this is insufficient.}

Throughout Sections~3.2, 3.4, and Appendix~F, every comparison between methods is presented as ``$\varepsilon$ achieved vs.\ computational cost.'' But $\varepsilon$ is a tuning parameter of the algorithm, not a measure of inferential quality. A method that reaches a smaller~$\varepsilon$ faster is not necessarily producing a better posterior approximation---it could be concentrating around a wrong mode, or suffering from particle degeneracy that is not reflected in~$\varepsilon$.

A JASA reader will expect at least one proper measure of posterior quality: marginal KL divergence against a reference posterior (available in closed form for the Gaussian toy model), Wasserstein distance, coverage of credible intervals, or squared error on point estimates. Without this, the central claim---that \textsc{permABC} methods are ``more efficient''---is only partially supported.

\major{2} \textbf{No comparison with standard ABC-SMC (without permutations).}

The paper compares \textsc{permABC-SMC} with Vanilla ABC, ABC-PMC, and ABC-Gibbs. But the natural baseline---standard ABC-SMC with the same adaptive $\varepsilon$ schedule of Del~Moral et~al.\ (2012), just without the permutation step---is absent. This makes it impossible to isolate the contribution of the permutation idea from the choice of SMC framework. The first question any reviewer will ask is: ``How much of the gain comes from permutations, and how much from switching to SMC?''

\major{3} \textbf{Theorem~1 is correct but limited, and the paper oversells it.}

Theorem~1 states that \textsc{permABC} $=$ ABC when $\varepsilon < \varepsilon^*$. This is clean but says little about the practically relevant regime where $\varepsilon \geq \varepsilon^*$. The paper claims ``the efficiency gains provided by \textsc{permABC} come at no cost to theoretical consistency'' (Section~1.4), but this is only true asymptotically as $\varepsilon \to 0$. For any fixed~$\varepsilon$ used in practice, the two posteriors may differ, and there is no bound on the discrepancy. The Corollary is a direct consequence and adds no new content. I would suggest either:
\begin{itemize}[nosep]
    \item providing a finite-$\varepsilon$ bound on the TV or Wasserstein distance between $\pi^*_\varepsilon$ and $\pi_\varepsilon$, or
    \item toning down the claims and clearly stating the limitation.
\end{itemize}

\major{4} \textbf{The Over-Sampling asymptotic results ($M \to \infty$) are stated informally.}

The new content in Section~2.3 (pushforward formulation, $M \to \infty$ convergence) is the most theoretically interesting part of the paper. However, the results are presented as displayed equations rather than as formal propositions with explicit assumptions and proofs. The key convergence claim ($d_M \to \delta_\beta$ a.s.) relies on ``mild regularity assumptions on $d$'' that are never stated. For JASA, these results should be either:
\begin{itemize}[nosep]
    \item formalized as proper propositions with proofs (even short ones), or
    \item clearly labeled as heuristic/conjectural with a pointer to where rigorous treatment will appear.
\end{itemize}

I note that the authors appear to have additional material with a more complete treatment including a $K \to \infty$ WABC connection. This would substantially strengthen the paper and should be integrated.

\major{5} \textbf{The SIR application is a proof of concept, not a validation.}

The SIR section is well-written and honest about model misspecification. However:
\begin{itemize}[nosep]
    \item There is no comparison with any other method on this problem (not even standard ABC-SMC). The reader cannot assess whether \textsc{permABC} is actually better here, or just the only thing that was tried.
    \item The noise model (log-normal on $\delta$ with $\sigma = 0.05$) is introduced but not justified. Why this form? Why this value of~$\sigma$? Was there a sensitivity analysis?
    \item The posterior on $\nu_k$ reaching values of~4 is acknowledged as implausible. This raises the question: is \textsc{permABC} producing a good approximation of a bad model, or a bad approximation? Without a simulated-data experiment where the ground truth is known, it is impossible to tell.
    \item The claim that ``\textsc{permABC} was specifically designed for these challenges'' positions the SIR as the motivating application, yet the numerical evidence is limited to two figures (posterior of $R_0$ and a map). At minimum, I would expect posterior predictive checks or a simulation study on synthetic SIR data with known parameters.
\end{itemize}

\section*{Minor concerns}

\minor{1} \textbf{Section~2 starts with a \texttt{\textbackslash subsection} and no introductory text.} The reader jumps from the section heading directly into ``State of the art.'' A brief paragraph outlining the three strategies would help.

\minor{2} \textbf{The Gaussian toy model (Section~3.2) never specifies the values of $(a, b, s, n, K)$.} The reader is told this is ``a classical hierarchical model'' but the exact experimental setup (prior hyperparameters, sample size per compartment, values of $K$ tested) is not given in the main text.

\minor{3} \textbf{Figure~2 (Over-Sampling posterior evolution) is never properly discussed.} The caption says ``Gaussian toy model'' without specifying which parameters. More importantly, the text around Figure~2 now contains the formal $M \to \infty$ analysis but the figure is not interpreted in light of that analysis. The prior-like behavior of $\beta$ for large $M$ and the concentration of the local parameters should be explicitly pointed out on the figure.

\minor{4} \textbf{The ABC-Gibbs comparison (Section~3.3) uses a deliberately adversarial model.} The ridge structure in the over-parameterized model is designed to make Gibbs fail. While the point is valid, a fairer comparison would also show a model where ABC-Gibbs works well (e.g.\ the standard Gaussian toy model of Section~3.2 with weak dependencies), to demonstrate that the methods have complementary strengths rather than presenting a one-sided comparison.

\minor{5} \textbf{Tracked-change markup (\texttt{\textbackslash old\{\}}, \texttt{\textbackslash new\{\}}) is still visible in the compiled PDF} (rendered as red strikethrough and blue text). These must be resolved before submission.

\minor{6} \textbf{Notation overload on $M$.} In Section~2.1, $M$ denotes the number of datasets simulated per parameter value. In Section~2.3 (Over-Sampling), $M$ denotes the number of simulated compartments. The reuse is potentially confusing and should be disambiguated.

\minor{7} \textbf{The ``unique particle rate'' (Section~2.1)} is proposed as a complementary diagnostic for sample diversity but is never reported in any experiment. Either use it or remove the discussion.

\minor{8} \textbf{Missing discussion of computational cost.} The $\mathcal{O}(K^3)$ cost of the assignment step is mentioned but never benchmarked. For $K = 94$ (SIR), what fraction of the total runtime is spent in the assignment step versus simulation? This matters for practitioners deciding whether to adopt the method.

\section*{Questions for the authors}

\begin{enumerate}[nosep]
    \item Can you provide posterior quality metrics (KL, $W_2$, coverage) for the Gaussian toy model where the posterior is available in closed form?
    \item What happens to \textsc{permABC} when exchangeability is only approximate (e.g.\ mildly heterogeneous priors across compartments)?
    \item For the SIR application, have you run a simulation study with known ground truth to validate the method before applying it to real data?
    \item The asymptotic material connecting \textsc{permABC} to WABC (compartment-level Wasserstein viewpoint, $K \to \infty$ limit) seems highly relevant. Why was it not included?
\end{enumerate}

\section*{Assessment}

The core idea of \textsc{permABC} is sound, practically motivated, and computationally elegant. The Over-Sampling and Under-Matching extensions are original contributions. However, the paper needs:
\begin{enumerate}[nosep]
    \item at least one proper posterior quality metric beyond $\varepsilon$,
    \item a comparison with standard ABC-SMC (without permutations),
    \item either formal statements of the $M \to \infty$ results or integration of the existing asymptotic material,
    \item a simulation-based validation of the SIR model (synthetic data with known truth).
\end{enumerate}

\noindent With these additions, the paper would make a solid contribution to JASA.

\end{document}
